\Box{22.00}{79.00}{14.00}{\digits\bfseries\fontsize{20}{25}\selectfont\color{blue}01}
\Box{40.00}{79.00}{148.00}{\sffamily\bfseries\fontsize{11.2}{16}\selectfont 全球协作的深度，塑造开放生态的体系能力}
\Box{40.00}{89.00}{148.00}{\fontsize{10.4}{16.4}\selectfont 欧洲的人才积累、中国的模型供给、印度与巴西的参与基础，体现全球开源的地域差异；增长速度与既有规模也不总是同向。将人才、组织、项目、代码、模型与数据统一观测，才能解释不同地区在创造、复用与维护中的分工及其时序变化。本报告主张把跨境接续、跨层支撑与长期维护纳入发展评价，以协作关系判断体系能力，为国际合作和公共投入提供依据。\textsuperscript{\color{teal}[1]}}
\Box{22.00}{125.00}{14.00}{\digits\bfseries\fontsize{20}{25}\selectfont\color{blue}02}
\Box{40.00}{125.00}{148.00}{\sffamily\bfseries\fontsize{11.2}{16}\selectfont 中国模型开源的领先，需要拓展为体系与协作优势}
\Box{40.00}{135.00}{148.00}{\fontsize{10.4}{16.4}\selectfont HF 期后观察中，Qwen 的衍生仓库达到 15.1 万个，显示以其为代表的中国模型处于全球开放复用生态的领先行列。围绕模型，训练推理、AI 工具链与异构算力适配正在分层展开；国家政策纳入开源贡献激励，xLLM 等项目进入基金会治理。下一步应以开放接口、共同评测与跨境维护连接国内产业和全球社区，使模型优势沉淀为长期体系能力。\textsuperscript{\color{teal}[2]}}
\Box{22.00}{171.00}{14.00}{\digits\bfseries\fontsize{20}{25}\selectfont\color{blue}03}
\Box{40.00}{171.00}{148.00}{\sffamily\bfseries\fontsize{11.2}{16}\selectfont 全球南方的潜力，在于成为共创者与治理参与者}
\Box{40.00}{181.00}{148.00}{\fontsize{10.4}{16.4}\selectfont 本地语言、文化知识与真实场景，是全球南方参与 AI 开源的独特资源。Masakhane 的非洲语言协作展示了以数据共建、模型适配和社区研究进入全球创新网络的路径。基金会应发挥连接长期资金、中立治理、人才培养与跨境合作的关键作用。AI 普惠应支持当地伙伴共同定义问题、掌握维护能力并参与规则制定，把资源可及推进为可持续公共能力。\textsuperscript{\color{teal}[3]}}
\Box{22.00}{217.00}{14.00}{\digits\bfseries\fontsize{20}{25}\selectfont\color{blue}04}
\Box{40.00}{217.00}{148.00}{\sffamily\bfseries\fontsize{11.2}{16}\selectfont 杭州等成长节点提示：中国 AI 开源需要因地施策}
\Box{40.00}{227.00}{148.00}{\fontsize{10.4}{16.4}\selectfont 中国 AI 开源呈现企业牵引、科研供给与社区接续并存的结构，模型、工具链与基础软件承担不同产业角色。人才规模、增长速度与模型供给指向不同的优势节点；杭州的开源人才增长与浙江的模型供给值得联读。地方决策应区分研究、工程适配、长期维护与应用人才需求，围绕自身产业补齐协作环节，通过跨城、跨境项目合作形成特色。\textsuperscript{\color{teal}[4]}}
